\documentclass[11pt,openany]{report}
\usepackage[letterpaper,margin=0.82in,headheight=14pt]{geometry}
\usepackage[T1]{fontenc}
\usepackage[utf8]{inputenc}
\usepackage{lmodern}
\usepackage{microtype}
\usepackage{amsmath,amssymb,amsthm,mathtools}
\usepackage{booktabs,longtable,tabularx,array,multirow}
\usepackage{graphicx,float,placeins}
\usepackage[dvipsnames]{xcolor}
\usepackage{enumitem}
\usepackage{fancyhdr}
\usepackage{listings}
\usepackage{url}
\usepackage[hidelinks]{hyperref}
\usepackage[nameinlink,noabbrev]{cleveref}
\usepackage{needspace}
\allowdisplaybreaks

\hypersetup{
  pdftitle={Projector-Generated Special Functions VII},
  pdfauthor={Komorebi P.},
  pdfsubject={SU(1,1) Racah packet recoupling, Casimir multiplicity spaces, pentagon coherence, and bistochastic quantum kernels}
}

\definecolor{Ink}{HTML}{17202A}
\definecolor{Blue}{HTML}{2E86AB}
\definecolor{Teal}{HTML}{2A9D8F}
\definecolor{Orange}{HTML}{F4A261}
\definecolor{Magenta}{HTML}{A23E8C}
\definecolor{Red}{HTML}{E76F51}
\definecolor{Paper}{HTML}{FAFAF7}
\definecolor{Muted}{HTML}{5D6D7E}
\definecolor{CodeBG}{HTML}{F3F5F5}

\pagecolor{Paper}
\color{Ink}
\graphicspath{{figures/}}
\setlength{\parindent}{0pt}
\setlength{\parskip}{0.52em}
\raggedbottom

\pagestyle{fancy}
\fancyhf{}
\fancyhead[L]{\small Projector-Generated Special Functions VII}
\fancyhead[R]{\small Komorebi P.}
\fancyfoot[C]{\thepage}
\renewcommand{\headrulewidth}{0.4pt}

\newtheorem{theorem}{Theorem}[chapter]
\newtheorem{proposition}[theorem]{Proposition}
\newtheorem{lemma}[theorem]{Lemma}
\newtheorem{corollary}[theorem]{Corollary}
\newtheorem{conjecture}[theorem]{Conjecture}
\theoremstyle{definition}
\newtheorem{definition}[theorem]{Definition}
\newtheorem{example}[theorem]{Worked example}
\newtheorem{algorithm}[theorem]{Algorithm}
\theoremstyle{remark}
\newtheorem{remark}[theorem]{Remark}
\newtheorem{boundary}[theorem]{Boundary statement}

\newcommand{\C}{\mathbb C}
\newcommand{\R}{\mathbb R}
\newcommand{\N}{\mathbb N}
\newcommand{\e}{\mathrm e}
\newcommand{\ii}{\mathrm i}
\newcommand{\dd}{\,\mathrm d}
\newcommand{\norm}[1]{\left\lVert#1\right\rVert}
\newcommand{\abs}[1]{\left\lvert#1\right\rvert}
\newcommand{\Expect}{\mathbb E}
\newcommand{\Var}{\operatorname{Var}}
\newcommand{\diag}{\operatorname{diag}}
\newcommand{\Id}{\mathrm I}
\newcommand{\Tr}{\operatorname{Tr}}
\newcommand{\rank}{\operatorname{rank}}
\newcommand{\spanop}{\operatorname{span}}
\newcommand{\Cas}{\mathcal C}
\newcommand{\Mcal}{\mathcal M}
\newcommand{\Rcal}{\mathcal R}
\newcommand{\Pcal}{\mathcal P}
\newcommand{\Hcal}{\mathcal H}
\newcommand{\ket}[1]{\left|#1\right\rangle}
\newcommand{\bra}[1]{\left\langle#1\right|}
\newcommand{\ip}[2]{\left\langle#1,#2\right\rangle}
\newcommand{\poch}[2]{\left(#1\right)_{#2}}
\newcommand{\Sh}{\operatorname{H}}

\newcolumntype{L}[1]{>{\raggedright\arraybackslash}p{#1}}

\lstdefinestyle{pythonstyle}{
  language=Python,
  basicstyle=\ttfamily\scriptsize,
  backgroundcolor=\color{CodeBG},
  frame=single,
  rulecolor=\color{Muted!45},
  breaklines=true,
  breakatwhitespace=false,
  showstringspaces=false,
  numbers=left,
  numberstyle=\tiny\color{Muted},
  numbersep=7pt,
  tabsize=4,
  columns=fullflexible,
  keepspaces=true
}
\lstset{style=pythonstyle}

\newcommand{\FigureWide}[3]{%
\Needspace{0.52\textheight}
\begin{figure}[H]
\centering
\includegraphics[width=0.92\textwidth,height=0.43\textheight,keepaspectratio]{#1.pdf}
\caption{#2 #3}
\label{fig:#1}
\end{figure}
}

\newcommand{\InlinePart}[1]{%
\FloatBarrier
\Needspace{10\baselineskip}
\refstepcounter{part}%
\addcontentsline{toc}{part}{\protect\numberline{\thepart}#1}%
\vspace{1.4em}
{\color{Blue}\Large\bfseries Part \thepart\par}
\vspace{0.2em}
{\color{Blue}\LARGE\bfseries #1\par}
\vspace{0.9em}
}

\newcommand{\FlowChapter}[1]{%
\FloatBarrier
\Needspace{14\baselineskip}
\refstepcounter{chapter}%
\setcounter{section}{0}%
\setcounter{subsection}{0}%
\addcontentsline{toc}{chapter}{\protect\numberline{\thechapter}#1}%
\markboth{#1}{#1}%
\vspace{1.3em}
{\color{Muted}\large\bfseries Chapter \thechapter\par}
\vspace{0.15em}
{\LARGE\bfseries #1\par}
\vspace{0.9em}
}

\title{\textbf{Projector-Generated Special Functions VII}\\[0.25em]
\Large $\mathrm{SU}(1,1)$ Racah Packet Recoupling, Casimir Multiplicity Spaces, Pentagon Coherence, and Bistochastic Quantum Kernels}
\author{Komorebi P.}
\date{}

\begin{document}
\hypersetup{pageanchor=false}
\pagenumbering{roman}
\maketitle

\begin{abstract}
The preceding paper produced a continuous shape semigroup, Meixner--Pollaczek packet bases, beta-binomial refinement, and a common positive discrete-series representation of $\mathfrak{su}(1,1)$. The present paper studies what happens when three or more shape packets are coupled. The result is a finite family of orthogonal recoupling functions generated by pair-Casimir projectors.

For positive shapes $\kappa_1,\kappa_2,\kappa_3$, let $D^+_{\kappa_j}$ denote the positive discrete-series representation. At fixed total excitation $N$, the tensor product contains an $(N+1)$-dimensional lowest-weight multiplicity space. Diagonalizing the intermediate Casimir $\Cas_{12}$ gives the left coupling atlas, while diagonalizing $\Cas_{23}$ gives the right coupling atlas. Their overlap is the Racah packet matrix
\[
\boxed{
\Rcal_N^{\boldsymbol\kappa}(j,m)
=
\ip{j,(12)3}{m,1(23)}.
}
\]
It is real orthogonal. In the $\Cas_{12}$ basis, the operator $\Cas_{23}$ is tridiagonal, so every column obeys a finite three-term recurrence. This is the packet-calculus realization of the classical Racah-polynomial and $6j$ structure.

Two exact low-degree transforms are derived. At $N=1$, the recoupling matrix is a shape-controlled beam splitter. At equal shape $\kappa_1=\kappa_2=\kappa_3=1/2$ and $N=2$, the complete $3\times3$ matrix is expressed through rational numbers and square roots. Squaring the matrix entries produces a doubly stochastic kernel. It acts as the transition matrix obtained by measuring one coupling observable after dephasing in the other coupling basis. The resulting channel preserves the uniform law and cannot decrease Shannon entropy.

For four packets, the fixed-excitation lowest-weight space has dimension $\binom{N+2}{2}$. Each binary coupling tree determines a complete joint eigenbasis of commuting intermediate Casimirs. The five trees of the associahedron are linked by orthogonal packet transforms. Their ordered product around the pentagon is exactly the identity. This gives a finite numerical and symbolic model of associator coherence within the projector calculus.

Every finite recoupling matrix can be compiled into at most $d(d-1)/2$ Givens rotations followed by diagonal signs. The full tensor space is needed only to define the lowest-weight projector. All later calculations take place in the multiplicity space, which gives an output-sensitive reduction.

Racah polynomials, $6j$ symbols, positive discrete-series recoupling, and the pentagon identity are classical. The generated contribution is the common packet architecture joining projector reduction, exact coupling atlases, stochastic kernels, entropy production, associator coherence, and finite quantum compilation. All figures were generated by the accompanying Python program. The release contains 30 vector figures, matching PNG files, 13 built-in checks, 20 independent tests, exact tables, and the complete Python source.
\end{abstract}

\chapter*{Status and claim boundary}
\addcontentsline{toc}{chapter}{Status and claim boundary}
The paper proves finite-dimensional statements about the positive discrete-series tensor products used below. It constructs the excitation spaces, lowest-weight projectors, pair-Casimir restrictions, orthogonal recoupling matrices, bistochastic kernels, four-body tree bases, pentagon closure, and Givens decompositions directly from matrices generated by the representation.

The historical identification of the recoupling entries with Racah polynomials and $6j$ symbols is classical. The classical theory also places Hahn polynomials at the two-body coupling stage and Wilson functions in continuous-series extensions. The present use of the word generated refers to the way these objects emerge from the packet calculus and not to a claim that Racah polynomials or $6j$ symbols are newly discovered here.

The stochastic channel obtained by squaring an orthogonal recoupling matrix describes basis change followed by dephasing or projective measurement. It does not represent coherent unitary evolution by itself. The entropy theorem applies to the diagonal probability vector after this measurement channel.

The pentagon product is exact because each tree basis is complete in the same multiplicity space. The numerical construction confirms the projector and ordering conventions. It is not offered as a new proof of the abstract monoidal coherence theorem.

The brief quantum discussion keeps the logarithmic Hamiltonian of the attached 2026 zeta-governed dynamics separate from the shape recoupling operator. The two modules can be composed, but their spectra are not identified. No result in this paper proves the Riemann hypothesis.

\begingroup
\footnotesize
\setlength{\parskip}{0pt}
\tableofcontents
\listoffigures
\endgroup
\clearpage
\pagenumbering{arabic}
\hypersetup{pageanchor=true}

\InlinePart{Positive discrete-series packet spaces}

\FlowChapter{From shape refinement to recoupling}

Paper VI produced the orthonormal Meixner--Pollaczek packet basis and the ladder
\[
K_0e_n=\left(n+\frac{\alpha}{2}\right)e_n,
\qquad
K_+e_n=\sqrt{(n+1)(n+\alpha)}\,e_{n+1},
\]
\[
K_-e_n=\sqrt{n(n+\alpha-1)}\,e_{n-1}.
\]
Writing
\[
\kappa=\frac{\alpha}{2}
\]
gives the standard positive discrete-series form.

\begin{definition}[Positive shape packet]
For $\kappa>0$, let $D^+_\kappa$ be the Hilbert space with orthonormal basis $\{\ket n:n\ge0\}$ and operators
\[
K_0\ket n=(n+\kappa)\ket n,
\]
\[
K_+\ket n=\sqrt{(n+1)(n+2\kappa)}\ket{n+1},
\]
\[
K_-\ket n=\sqrt{n(n+2\kappa-1)}\ket{n-1}.
\]
\end{definition}

Direct calculation gives
\[
[K_0,K_\pm]=\pm K_\pm,
\qquad
[K_-,K_+]=2K_0.
\]
The Casimir
\[
\boxed{
\Cas=K_0(K_0-1)-K_+K_-
}
\]
acts as the scalar
\[
\boxed{
\Cas=\kappa(\kappa-1)\Id.
}
\]

For $m$ packets, the total generators are obtained from the primitive coproduct
\[
K_\bullet^{[m]}=\sum_{r=1}^{m}K_\bullet^{(r)}.
\]
The operator $K_0^{[m]}$ resolves the tensor product by total excitation.

\FigureWide{01_theory_pipeline}{From shape refinement to Racah recoupling.}{The same positive shape packet that generated the Meixner basis becomes the input to a tensor-product projector, a pair-Casimir atlas, a Racah associator, and a finite stochastic or quantum transform.}

\FlowChapter{Fixed excitation and the lowest-weight projector}

Let
\[
\boldsymbol\kappa=(\kappa_1,\ldots,\kappa_m)
\]
and define the total-excitation space
\[
\Hcal_N^{(m)}
=
\spanop\left\{
\ket{n_1,\ldots,n_m}
:
\sum_{r=1}^{m}n_r=N
\right\}.
\]
Its dimension is the number of weak compositions of $N$ into $m$ parts.

\begin{proposition}[Excitation-space dimension]
\[
\boxed{
\dim\Hcal_N^{(m)}
=
\binom{N+m-1}{m-1}.
}
\]
\end{proposition}

The total lowering operator maps
\[
K_-^{[m]}:\Hcal_N^{(m)}\longrightarrow\Hcal_{N-1}^{(m)}.
\]
The packet of new lowest-weight vectors appearing at level $N$ is
\[
\boxed{
\Mcal_N^{(m)}
=
\ker\left(K_-^{[m]}\big|_{\Hcal_N^{(m)}}\right).
}
\]

\begin{theorem}[Multiplicity dimension]
For $N\ge0$ and $m\ge2$,
\[
\boxed{
\dim\Mcal_N^{(m)}
=
\binom{N+m-2}{m-2}.
}
\]
\end{theorem}

\begin{proof}
The adjoint of the lowering map is the total raising map
\[
K_+^{[m]}:\Hcal_{N-1}^{(m)}\longrightarrow\Hcal_N^{(m)}.
\]
The positive discrete series has no highest-weight vector, so $K_+^{[m]}$ is injective on every finite excitation level. Hence the lowering map has full row rank. Rank-nullity gives
\[
\begin{aligned}
\dim\ker K_-^{[m]}
&=
\binom{N+m-1}{m-1}
-
\binom{N+m-2}{m-1}\\
&=
\binom{N+m-2}{m-2}.
\end{aligned}
\]
\end{proof}

The orthogonal projector onto $\Mcal_N^{(m)}$ is the first major compression in the construction. The ratio between the full excitation dimension and the multiplicity dimension is
\[
\boxed{
\frac{\dim\Hcal_N^{(m)}}{\dim\Mcal_N^{(m)}}
=
\frac{N+m-1}{m-1}.
}
\]

\FigureWide{02_su11_ladder}{Positive discrete-series ladder.}{Each shape parameter gives an infinite ladder. Tensor products are first resolved by the total excitation and then projected onto the new lowest-weight multiplicity packet.}

\FigureWide{03_excitation_simplex}{The total-excitation simplex.}{Every lattice point is a composition of the total degree. The color records one occupation number and makes the finite packet geometry visible.}

\FigureWide{04_dimension_compression}{Projector reduction before recoupling.}{The full tensor excitation space grows faster than the lowest-weight multiplicity space. The later Racah calculation uses only the multiplicity packet.}

\FlowChapter{Intermediate Casimir projectors}

For a subset $A\subseteq\{1,\ldots,m\}$, define
\[
K_\bullet^{A}=\sum_{r\in A}K_\bullet^{(r)}
\]
and
\[
\boxed{
\Cas_A
=
K_0^A(K_0^A-1)-K_+^AK_-^A.
}
\]
If two subsets are nested or disjoint, the corresponding Casimirs commute.

For two packets,
\[
D^+_{\kappa_a}\otimes D^+_{\kappa_b}
\cong
\bigoplus_{j=0}^{\infty}
D^+_{\kappa_a+\kappa_b+j}.
\]
Therefore the intermediate Casimir spectrum at total excitation $N$ is
\[
\boxed{
\lambda_j^{ab}
=
(\kappa_a+\kappa_b+j)
(\kappa_a+\kappa_b+j-1),
\qquad
0\le j\le N.
}
\]

\begin{proposition}[Three-packet pair spectrum]
On $\Mcal_N^{(3)}$, both $\Cas_{12}$ and $\Cas_{23}$ have simple spectra of length $N+1$. Their eigenvalues are $\lambda_j^{12}$ and $\lambda_m^{23}$ respectively.
\end{proposition}

\begin{proof}
The two-packet decomposition above fixes the intermediate shape. Coupling the resulting representation to the remaining packet gives the final lowest weight
\[
\kappa_1+\kappa_2+\kappa_3+N.
\]
For every $j\in\{0,\ldots,N\}$ there is exactly one compatible remaining excitation $N-j$. Thus each intermediate shape appears once in the fixed final multiplicity space.
\end{proof}

\FigureWide{05_pair_casimir_spectra}{Pair-Casimir spectra.}{Both coupling schemes have simple quadratic spectra. Their eigenvectors, not their eigenvalues, carry the nontrivial recoupling information.}

\InlinePart{Three-body Racah packet functions}

\FlowChapter{Two complete coupling atlases}

For three packets, define the left basis by
\[
\Cas_{12}\ket{j,(12)3}
=
\lambda_j^{12}\ket{j,(12)3}
\]
and the right basis by
\[
\Cas_{23}\ket{m,1(23)}
=
\lambda_m^{23}\ket{m,1(23)}.
\]
Both bases are complete in $\Mcal_N^{(3)}$.

\begin{definition}[Racah packet matrix]
\[
\boxed{
\Rcal_N^{\boldsymbol\kappa}(j,m)
=
\bra{j,(12)3}\ket{m,1(23)}.
}
\]
\end{definition}

After a real phase convention, the matrix $\Rcal_N^{\boldsymbol\kappa}$ is real orthogonal.

\begin{theorem}[Perfect recoupling reconstruction]
\[
\boxed{
\sum_{m=0}^{N}
\Rcal_N^{\boldsymbol\kappa}(j,m)
\Rcal_N^{\boldsymbol\kappa}(j',m)
=
\delta_{jj'}
}
\]
and
\[
\boxed{
\sum_{j=0}^{N}
\Rcal_N^{\boldsymbol\kappa}(j,m)
\Rcal_N^{\boldsymbol\kappa}(j,m')
=
\delta_{mm'}.
}
\]
\end{theorem}

\begin{proof}
These are the row and column orthogonality relations of the overlap between two orthonormal bases of the same finite-dimensional space.
\end{proof}

The entries are the positive discrete-series $6j$ coefficients, after normalization and phase choice. Their scalar polynomial realization is the Racah family. We retain the classical name and use packet to emphasize the analysis-synthesis role of the entire matrix.

\FlowChapter{One excitation as a shape beam splitter}

At $N=1$, the multiplicity space has dimension two. Write
\[
K=\kappa_1+\kappa_2+\kappa_3.
\]
The complete transform is
\[
\boxed{
\Rcal_1^{\boldsymbol\kappa}
=
\begin{pmatrix}
\displaystyle
\sqrt{\frac{\kappa_1\kappa_3}
{(\kappa_1+\kappa_2)(\kappa_2+\kappa_3)}}
&
\displaystyle
\sqrt{\frac{\kappa_2K}
{(\kappa_1+\kappa_2)(\kappa_2+\kappa_3)}}
\\[1.2em]
\displaystyle
\sqrt{\frac{\kappa_2K}
{(\kappa_1+\kappa_2)(\kappa_2+\kappa_3)}}
&
\displaystyle
-\sqrt{\frac{\kappa_1\kappa_3}
{(\kappa_1+\kappa_2)(\kappa_2+\kappa_3)}}
\end{pmatrix}.
}
\]

\begin{proof}
Let the two entries in the first row be $c$ and $s$. Their squares add to
\[
\frac{\kappa_1\kappa_3+\kappa_2K}
{(\kappa_1+\kappa_2)(\kappa_2+\kappa_3)}
=1.
\]
The lowest-weight vectors associated with the two intermediate Casimir eigenvalues determine the positive choice for the first row. Orthogonality then fixes the second row up to one sign. The displayed convention agrees with the ordered eigenbases used throughout the paper.
\end{proof}

This transform behaves like a shape-controlled beam splitter. Its transition probabilities are
\[
P=
\begin{pmatrix}
c^2&s^2\\
s^2&c^2
\end{pmatrix}.
\]
When $\kappa_1=\kappa_2=\kappa_3=1/2$,
\[
c=\frac12,
\qquad
s=\frac{\sqrt3}{2}.
\]

\FigureWide{06_shape_mixing_angle}{One-excitation shape mixing.}{The middle packet controls how strongly the two coupling atlases mix. The squared entries form a two-state bistochastic channel.}

\FigureWide{07_one_excitation_matrix}{Exact one-excitation matrix.}{The displayed numerical example agrees with the closed shape formula to floating-point roundoff.}

\FlowChapter{An exact three-channel example}

Take
\[
\kappa_1=\kappa_2=\kappa_3=\frac12
\]
and $N=2$. Then
\[
\boxed{
\Rcal_2
=
\begin{pmatrix}
\frac13&\frac1{\sqrt3}&\frac{\sqrt5}{3}\\[0.5em]
\frac1{\sqrt3}&\frac12&-\frac{\sqrt{15}}6\\[0.5em]
\frac{\sqrt5}{3}&-\frac{\sqrt{15}}6&\frac16
\end{pmatrix}.
}
\]

\begin{example}[Direct orthogonality check]
The first row has squared norm
\[
\frac19+\frac13+\frac59=1.
\]
The inner product of the first two rows is
\[
\frac1{3\sqrt3}
+
\frac1{2\sqrt3}
-
\frac{\sqrt{75}}{18}
=0.
\]
The remaining row relations follow in the same way.
\end{example}

The corresponding probability kernel is
\[
P=
\begin{pmatrix}
\frac19&\frac13&\frac59\\[0.3em]
\frac13&\frac14&\frac5{12}\\[0.3em]
\frac59&\frac5{12}&\frac1{36}
\end{pmatrix}.
\]
Every row and every column sums to one.

\FigureWide{08_two_excitation_exact}{Equal-shape recoupling at total excitation two.}{The matrix contains only rational numbers and square roots. Its signed entries are amplitudes and their squares are transition probabilities.}

\FlowChapter{The Leonard pair and the Racah recurrence}

Express $\Cas_{23}$ in the $\Cas_{12}$ eigenbasis:
\[
J^{(23|12)}
=
\Rcal_N
\diag(\lambda_0^{23},\ldots,\lambda_N^{23})
\Rcal_N^{\mathsf T}.
\]
The Racah algebra implies that this matrix is irreducibly tridiagonal.

\begin{theorem}[Finite three-term packet recurrence]
There are real coefficients $a_j$ and positive off-diagonal coefficients $b_j$ such that every column satisfies
\[
\boxed{
b_{j-1}\Rcal_N(j-1,m)
+a_j\Rcal_N(j,m)
+b_j\Rcal_N(j+1,m)
=
\lambda_m^{23}\Rcal_N(j,m).
}
\]
The endpoint terms outside $0\le j\le N$ are omitted.
\end{theorem}

\begin{proof}
The matrix $J^{(23|12)}$ is the representation of $\Cas_{23}$ in the $\Cas_{12}$ basis. The pair $\Cas_{12},\Cas_{23}$ generates the rank-one Racah algebra. On the irreducible fixed-$N$ multiplicity module, either Casimir is diagonal in its own basis and tridiagonal in the eigenbasis of the other. Applying the matrix eigenvalue equation to the $m$th column gives the recurrence.
\end{proof}

The dual statement holds with $j$ and $m$ exchanged. Thus the packet matrix is bispectral. Classical normalization turns these finite functions into Racah polynomials.

\FigureWide{09_racah_packet_heatmap}{Racah packet transform at total excitation twelve.}{The alternating bands show the finite oscillatory structure of the associator. Orthogonality holds across the complete matrix, not packet by packet.}

\FigureWide{12_leonard_tridiagonal}{Leonard-pair selection rule.}{In the left Casimir basis, the right Casimir connects only neighboring packet labels. This produces the finite three-term Racah recurrence.}

\FigureWide{13_racah_columns}{Racah packet columns.}{Selected columns behave as finite orthogonal special functions on the intermediate-Casimir lattice.}

\FlowChapter{Parameter dependence and localization}

The shapes $\kappa_1,\kappa_2,\kappa_3$ are continuous. The pair spectra and recoupling matrix vary smoothly as long as all shapes remain positive.

The maximal entry
\[
\mu_N(\boldsymbol\kappa)
=
\max_{j,m}\abs{\Rcal_N^{\boldsymbol\kappa}(j,m)}
\]
measures the coherence between the two coupling atlases. Small $\mu_N$ means that no basis vector is strongly concentrated in the other atlas.

\FigureWide{14_maximum_coherence}{Coupling-basis coherence across packet depth.}{The maximal amplitude decreases with dimension in the displayed shape family, while the equal-shape and unequal-shape cases follow different finite-depth profiles.}

\FigureWide{16_parameter_sensitivity}{Shape dependence of the largest recoupling amplitude.}{The heat map records how the most concentrated overlap changes when one shape is varied. It is a diagnostic rather than a universal monotonicity theorem.}

\InlinePart{Bistochastic and quantum packet channels}

\FlowChapter{From amplitudes to a Markov kernel}

Define
\[
\boxed{
P_{jm}
=
\abs{\Rcal_N(j,m)}^2.
}
\]

\begin{theorem}[Bistochastic Racah kernel]
The matrix $P$ is doubly stochastic:
\[
P_{jm}\ge0,
\qquad
\sum_mP_{jm}=1,
\qquad
\sum_jP_{jm}=1.
\]
\end{theorem}

\begin{proof}
Nonnegativity is immediate. Row sums and column sums are the squared norms of the rows and columns of the orthogonal matrix $\Rcal_N$.
\end{proof}

The channel has a direct measurement interpretation. Begin with a density matrix diagonal in the left coupling basis,
\[
\rho=\sum_jp_j\ket{j,(12)3}\bra{j,(12)3}.
\]
Measure in the right coupling basis and discard the outcome phase. The output distribution is
\[
\boxed{
q_m=\sum_jp_jP_{jm}.
}
\]

\FigureWide{10_bistochastic_kernel}{Squared recoupling amplitudes.}{The kernel is nonnegative and every row and column sums to one. It is the classical channel induced by coherent basis change followed by measurement.}

\FigureWide{17_bistochastic_marginals}{Exact row and column marginals.}{The numerical residuals remain at floating-point scale.}

\FlowChapter{Entropy production and uncertainty}

Let
\[
\Sh(p)=-\sum_jp_j\log p_j
\]
be the Shannon entropy.

\begin{theorem}[Entropy monotonicity]
For every probability vector $p$,
\[
\boxed{
\Sh(pP)\ge\Sh(p).
}
\]
\end{theorem}

\begin{proof}
A doubly stochastic matrix is a convex combination of permutation matrices. Therefore $pP$ is majorized by $p$. Shannon entropy is Schur concave, which gives the result.
\end{proof}

The uniform distribution is stationary:
\[
\frac1{N+1}(1,\ldots,1)P
=
\frac1{N+1}(1,\ldots,1).
\]

For a left basis state $j$, the right-basis entropy is the row entropy
\[
\Sh_j=-\sum_mP_{jm}\log P_{jm}.
\]
The Maassen--Uffink bound gives
\[
\Sh_{12}+\Sh_{23}
\ge
-2\log\mu_N.
\]
For a pure left basis packet, $\Sh_{12}=0$.

\FigureWide{15_row_entropy}{A basis packet spreads under recoupling.}{The solid curve is the mean row entropy and the band records the range across packet labels. The dotted line is the uniform maximum.}

\FlowChapter{Repeated measurement and mixing}

Repeated application of $P$ defines a finite Markov chain. Because $P$ is bistochastic, the uniform law is always stationary. When the kernel is primitive, every initial probability vector converges to uniform.

The coherent matrix $\Rcal_N$ and the classical kernel $P$ must remain distinct. Two coherent recouplings can interfere, while two applications of $P$ cannot recover phase.

\begin{boundary}[What entropy does not prove]
Entropy increase under $P$ does not mean that the unitary Racah transform is irreversible. The entropy appears only after dephasing or measurement discards the phase relations. The underlying orthogonal transform has the exact inverse $\Rcal_N^{\mathsf T}$.
\end{boundary}

\InlinePart{Four-body packet trees and pentagon coherence}

\FlowChapter{Five complete coupling trees}

For four packets and total excitation $N$,
\[
\dim\Mcal_N^{(4)}
=
\binom{N+2}{2}.
\]
A binary coupling tree is specified by two commuting intermediate Casimirs. The five trees are
\[
(((12)3)4),
\quad
((1(23))4),
\quad
(1((23)4)),
\quad
(1(2(34))),
\quad
((12)(34)).
\]
Each tree gives a complete joint eigenbasis of the same multiplicity space.

\begin{proposition}[Commuting tree observables]
For every tree, its two intermediate Casimirs commute on the complete tensor product and on the lowest-weight multiplicity packet.
\end{proposition}

\begin{proof}
The subsets labeling the two internal edges are nested or disjoint. Coproduct generators on disjoint tensor factors commute. A Casimir of a smaller subset commutes with the total generators of that subset and hence with the Casimir of a containing subset.
\end{proof}

\FigureWide{18_three_body_trees}{Two three-packet coupling trees.}{The Racah matrix is the exact basis transform between the two complete intermediate-Casimir atlases.}

\FigureWide{19_pentagon_trees}{The four-packet associahedron.}{The five binary trees form a pentagon. Each edge carries an orthogonal recoupling matrix.}

\FigureWide{22_joint_tree_spectrum}{Joint packet labels for one four-body tree.}{The two diagonal Casimirs provide a complete finite address on the multiplicity packet.}

\FlowChapter{The finite pentagon identity}

Let $B_0,\ldots,B_4$ be the five orthogonal tree-basis matrices, with columns expressed in one common multiplicity frame. Define the edge transform
\[
U_{r,r+1}=B_r^{\mathsf T}B_{r+1},
\qquad
r\pmod5.
\]

\begin{theorem}[Pentagon closure]
\[
\boxed{
U_{01}U_{12}U_{23}U_{34}U_{40}=\Id.
}
\]
\end{theorem}

\begin{proof}
Insert the definitions and use $B_rB_r^{\mathsf T}$ as the identity on the common multiplicity space:
\[
\begin{aligned}
U_{01}U_{12}U_{23}U_{34}U_{40}
&=
B_0^{\mathsf T}B_1
B_1^{\mathsf T}B_2
B_2^{\mathsf T}B_3
B_3^{\mathsf T}B_4
B_4^{\mathsf T}B_0\\
&=
B_0^{\mathsf T}B_0\\
&=
\Id.
\end{aligned}
\]
\end{proof}

This is the finite matrix form of associator coherence. It also acts as a stringent implementation certificate. A wrong phase, label order, or missing multiplicity vector breaks the product.

\FigureWide{20_pentagon_residual}{Pentagon closure on growing multiplicity spaces.}{The maximum residual remains near floating-point roundoff while the dimension increases.}

\FigureWide{21_pentagon_edge_kernels}{Five edge associators.}{Each panel is an orthogonal transform between neighboring coupling trees at the same total excitation.}

\FlowChapter{Coassociativity and packet reconstruction}

The beta-binomial refinement of Paper VI already proved that splitting a total shape into several child shapes is coassociative. The present paper adds the change of intermediate basis between different bracketings.

At the amplitude level, a state may be expanded in any tree basis:
\[
\ket\Psi
=
\sum_{\alpha}c_\alpha^{(T)}\ket{\alpha,T}.
\]
For neighboring trees,
\[
c^{(T')}=U_{T,T'}^{\mathsf T}c^{(T)}.
\]
Because the pentagon closes, the coefficient vector is independent of the chosen route through the coupling graph.

This is the packet form of path-independent reconstruction:
\[
\boxed{
\text{analysis in one tree}
\longrightarrow
\text{any legal sequence of recouplings}
\longrightarrow
\text{same synthesized state}.
}
\]

\FigureWide{23_four_body_compression}{Four-body projector reduction.}{The multiplicity packet is substantially smaller than the full excitation simplex. All five tree transforms live in the reduced space.}

\InlinePart{Finite compilation and algorithms}

\FlowChapter{Givens synthesis of a recoupling transform}

Any real orthogonal $d\times d$ matrix can be triangularized by plane rotations. Reversing the rotations gives a finite circuit.

\begin{theorem}[Givens packet compilation]
Every Racah packet matrix $R\in O(d)$ can be written as
\[
\boxed{
R=G_1^{\mathsf T}\cdots G_L^{\mathsf T}D,
}
\]
where $D$ is diagonal with entries $\pm1$, every $G_\ell$ acts on two coordinates, and
\[
\boxed{
L\le\frac{d(d-1)}2.
}
\]
\end{theorem}

\begin{proof}
Apply standard Givens elimination below the diagonal. At most one rotation is required for every strictly lower-triangular entry. Orthogonality forces the resulting upper-triangular matrix to be diagonal with signs. Reverse the elimination sequence.
\end{proof}

For the $8\times8$ example in the release, the compiler uses 28 rotations, which meets the triangular bound exactly.

\FigureWide{24_givens_network}{Triangular Givens network.}{Each vertical segment is a two-coordinate rotation. The diagonal signs finish the exact finite synthesis.}

\FlowChapter{Output-sensitive recoupling algorithm}

The computational route is:

\begin{algorithm}[Projector-resolved Racah transform]
Given positive shapes $\boldsymbol\kappa$ and total excitation $N$:
\begin{enumerate}[label=\arabic*.]
\item Enumerate the weak compositions of $N$.
\item Build the total lowering operator from the ladder coefficients.
\item Compute an orthonormal frame of its kernel.
\item Restrict the required pair Casimirs to that frame.
\item Diagonalize one or more commuting Casimirs for each coupling tree.
\item Form overlap matrices between the resulting bases.
\item Use the signed overlaps for coherent synthesis or square them for measurement channels.
\item Compile any selected overlap into Givens rotations when a finite circuit is required.
\end{enumerate}
\end{algorithm}

The full tensor level is needed only in steps 1 through 3. A symbolic or recursive implementation can avoid even that materialization by constructing Clebsch--Gordan packets directly.

\FigureWide{26_packet_algorithm}{Output-sensitive recoupling algorithm.}{Project first, then diagonalize and recouple inside the multiplicity space. Dense full-tensor operators are not preserved as an output requirement.}

\FlowChapter{Scaling and numerical conditioning}

For three packets, the multiplicity dimension is $N+1$. For four packets, it is $\binom{N+2}{2}$. Generic dense eigendecomposition in the reduced space scales cubically in that dimension.

The classical three-term recurrence provides a faster route for large $N$. Once the Jacobi coefficients are known, selected Racah columns can be evaluated in linear time. Fast transforms and structured factorizations remain an open algorithmic direction.

\FigureWide{25_runtime_scaling}{Reference complexity profiles.}{The deterministic curves compare a dense reduced eigensolve, a structured full transform, and one recurrence-generated column after normalization at the smallest displayed dimension.}

The dimension reduction shown in \cref{fig:04_dimension_compression} is the first source of computational savings.

\InlinePart{Classical identification, quantum use, and future deformations}

\FlowChapter{Classical Racah and Wilson coordinates}

The polynomial hierarchy is classical:
\[
\text{Meixner--Pollaczek}
\longrightarrow
\text{Hahn coupling coefficients}
\longrightarrow
\text{Racah recoupling coefficients}.
\]
Racah polynomials are finite ${}_4F_3$ functions. In the standard notation,
\[
R_n(\lambda(x)\mid\alpha,\beta,\gamma,\delta)
=
{}_4F_3\!\left(
\begin{matrix}
-n,n+\alpha+\beta+1,-x,x+\gamma+\delta+1\\
\alpha+1,\beta+\delta+1,\gamma+1
\end{matrix}
\,\middle|\,1
\right),
\]
with
\[
\lambda(x)=x(x+\gamma+\delta+1)
\]
and a terminating parameter that restricts the support.

The representation shapes determine the Racah parameters after the conventional normalization of the $6j$ symbols. We define the functions intrinsically through Casimir projectors because that form extends directly to multiple coupling trees and numerical packet construction.

Continuous-series recoupling replaces finite Racah polynomials by Wilson polynomials and Wilson functions. A $q$-deformation leads to $q$-Racah and Askey--Wilson structures. Higher tensor products produce multivariate Racah functions and general $3nj$ coefficients.

\FigureWide{27_classical_family_map}{Classical-family map.}{The generated packet language retains the classical names and adds a common analysis, synthesis, refinement, and associator interpretation.}

\FlowChapter{Relation to the logarithmic quantum coordinate}

The attached 2026 zeta-governed quantum construction uses the number basis and the Hamiltonian
\[
H\ket n=\log n\ket n.
\]
The resulting phase is
\[
\e^{-\ii tH}\ket n=n^{-\ii t}\ket n.
\]

The shape module in the present paper is different. Its basis labels the positive discrete-series packet degree and its finite recoupling matrix changes the bracketing of shape factors. A composite model may therefore use
\[
\boxed{
\text{logarithmic phase evolution}
\otimes
\text{finite shape recoupling}.
}
\]
The two operators commute when they act on separate registers. No claim is made that the zeta zeros are eigenvalues of the Racah transform.

\FigureWide{28_quantum_connection}{Two quantum coordinates.}{The logarithmic Hamiltonian acts on the number register, while the Racah transform acts on a finite shape-coupling register. Their roles remain distinct.}

\FlowChapter{What the packet calculus adds}

The classical recoupling coefficient becomes a more general coordinate object when placed inside the calculus. The complete data are
\[
\boxed{
\left(
\Mcal_N,
\{\Cas_A\},
\{B_T\},
\{U_{T,T'}\},
\{|U_{T,T'}|^2\},
\text{refinement graph},
\text{certificate}
\right).
}
\]
This package supports several outputs from one computation:

\begin{itemize}
\item coherent basis transformation
\item a classical measurement channel
\item entropy and uncertainty bounds
\item route-independent tree recoupling
\item finite circuit compilation
\item reduced-space numerical algorithms
\end{itemize}

The recoupling entries are not renamed as new scalar functions. The new object is the typed packet system that keeps all these roles compatible.

\FigureWide{29_theorem_map}{Theorem dependency map.}{Pair-Casimir projectors generate the orthogonal associator. Its signed and squared forms then support Racah recurrences, stochastic channels, pentagon coherence, and finite compilation.}

\FlowChapter{Open problems generated by the construction}

The following questions are concrete continuations.

\begin{enumerate}
\item Develop fast transforms for large finite Racah packet matrices with provable complexity below dense quadratic application.
\item Derive stable packet recursion across both excitation $N$ and continuous shapes $\boldsymbol\kappa$.
\item Establish sharp coherence and entropy bounds uniform in the packet labels.
\item Construct multivariate Racah packet bases directly from the full tree rather than repeated pair transforms.
\item Extend the projector construction to the principal and complementary series, where Wilson functions replace finite Racah polynomials.
\item Develop the $q$-deformed packet calculus and compare the resulting associator with quantum-group circuit models.
\item Connect packet resultants and Casimir spectra to exact denominator and arithmetic properties at rational shape parameters.
\end{enumerate}

\FigureWide{30_future_directions}{Next deformation directions.}{Continuous, $q$-deformed, multivariate, and quantum-circuit extensions retain the same basic rule: declare a complete coupling atlas, compute the associator, and certify reconstruction.}

\FlowChapter{Conclusion}

The shape parameter introduced in Paper VI carries more structure than a convolution coordinate. Its positive discrete-series representation can be coupled, projected, and recoupled. At fixed total excitation, the new lowest-weight multiplicity space is finite. Pair Casimirs provide complete coordinate atlases, and their overlap is a Racah packet transform.

The signed transform is orthogonal. Its squared entries form a doubly stochastic channel. Four-packet coupling trees satisfy an exact pentagon closure. Every finite transform admits a Givens circuit. These facts arise from one projector-resolved object rather than from separate calculations.

The classical Racah and $6j$ theory supplies the scalar special functions. The calculus supplies the common language of packets, projectors, stochastic shadows, tree refinement, and exact reconstruction.

\appendix

\chapter{Formula registry}

\begin{longtable}{L{0.30\textwidth}L{0.60\textwidth}}
\toprule
Object & Exact formula or property\\
\midrule
\endfirsthead
\toprule
Object & Exact formula or property\\
\midrule
\endhead
Positive discrete-series Casimir & $\Cas=\kappa(\kappa-1)\Id$\\
Excitation dimension & $\binom{N+m-1}{m-1}$\\
Multiplicity dimension & $\binom{N+m-2}{m-2}$\\
Pair-Casimir eigenvalue & $(\kappa_a+\kappa_b+j)(\kappa_a+\kappa_b+j-1)$\\
Racah packet & $\Rcal_N(j,m)=\ip{j,(12)3}{m,1(23)}$\\
Orthogonality & $\Rcal_N\Rcal_N^{\mathsf T}=\Id$\\
Bistochastic kernel & $P_{jm}=|\Rcal_N(j,m)|^2$\\
Entropy law & $\Sh(pP)\ge\Sh(p)$\\
Four-body multiplicity & $\binom{N+2}{2}$\\
Pentagon product & $U_{01}U_{12}U_{23}U_{34}U_{40}=\Id$\\
Givens count & at most $d(d-1)/2$ rotations\\
\bottomrule
\end{longtable}

\chapter{Numerical certificate tables}

\section{Three-body metrics}
\begin{center}
\small
\begin{tabular}{rrrrr}
\toprule
$N$ & dimension & largest $|R_{jm}|$ & mean row entropy & tridiagonal residual\\
\midrule
1&2&0.90830&0.46373&0\\
2&3&0.81922&0.86600&$1.0\times10^{-16}$\\
4&5&0.68308&1.20209&$1.5\times10^{-14}$\\
7&8&0.63965&1.62054&$4.0\times10^{-14}$\\
12&13&0.57135&2.05931&$1.2\times10^{-13}$\\
\bottomrule
\end{tabular}
\end{center}

\section{Four-body pentagon metrics}
\begin{center}
\small
\begin{tabular}{rrrr}
\toprule
$N$ & multiplicity dimension & pentagon residual & largest joint-diagonalization residual\\
\midrule
1&3&$1.44\times10^{-15}$&$4.41\times10^{-15}$\\
2&6&$1.67\times10^{-15}$&$1.18\times10^{-14}$\\
3&10&$3.11\times10^{-15}$&$1.61\times10^{-13}$\\
4&15&$3.44\times10^{-15}$&$6.58\times10^{-13}$\\
7&36&$3.49\times10^{-15}$&$2.65\times10^{-12}$\\
\bottomrule
\end{tabular}
\end{center}

\chapter{Verification ledger}

The built-in program checks the following statements.

\begin{longtable}{L{0.50\textwidth}L{0.18\textwidth}L{0.22\textwidth}}
\toprule
Certificate & Result & Largest recorded residual\\
\midrule
\endfirsthead
\toprule
Certificate & Result & Largest recorded residual\\
\midrule
\endhead
Pair-Casimir spectra & pass & $2.84\times10^{-14}$\\
Three-body orthogonality & pass & $1.33\times10^{-15}$\\
One-excitation closed form & pass & $2.78\times10^{-16}$\\
Two-excitation exact matrix & pass & $6.11\times10^{-16}$\\
Racah tridiagonal selection & pass & $3.97\times10^{-14}$\\
Bistochastic kernel & pass & $1.33\times10^{-15}$\\
Entropy monotonicity & pass & minimum gap $2.18\times10^{-2}$\\
Multiplicity dimension & pass & exact\\
Pentagon coherence & pass & $3.44\times10^{-15}$\\
Tree Casimir commutation & pass & $2.84\times10^{-13}$\\
Joint diagonalization & pass & $6.58\times10^{-13}$\\
Givens reconstruction & pass & $9.99\times10^{-16}$\\
Equal-shape probability symmetry & pass & $1.33\times10^{-15}$\\
\bottomrule
\end{longtable}

The independent test file contains 20 additional tests. It rechecks dimensions, nullspace relations, several shape choices, exact low-degree formulas, bistochasticity, entropy, pentagon closure, orthogonality of every edge transform, and source artifacts.

\chapter{Complete Python implementation}

The following file generates every figure, table, built-in certificate, and release manifest used in the paper.

\lstinputlisting{projector_generated_special_functions_vii.py}

\chapter{Bibliographic note}

\begin{thebibliography}{10}

\bibitem{DLMFRacah}
NIST Digital Library of Mathematical Functions,
Sections 18.19, 18.20, 18.21, 18.25, and 18.38,
Hahn, Racah, Wilson, and $6j$ relations.
\url{https://dlmf.nist.gov/18}

\bibitem{KoelinkVdJeugt}
H. T. Koelink and J. Van der Jeugt,
Convolutions for orthogonal polynomials from Lie and quantum algebra representations,
\emph{SIAM Journal on Mathematical Analysis} 29 (1998), 794--822.

\bibitem{Groenevelt}
W. Groenevelt,
Wilson function transforms related to Racah coefficients,
\emph{Acta Applicandae Mathematicae} 91 (2006), 133--191.

\bibitem{LievensVdJeugt}
S. Lievens and J. Van der Jeugt,
Racah coefficients of the positive discrete series representations of $\mathfrak{su}(1,1)$,
representation-theoretic recoupling literature.

\bibitem{Genest}
V. X. Genest, L. Vinet, and collaborators,
Multivariate Racah and $q$-Racah polynomials as coupling coefficients,
works on higher-rank recoupling and quantum-algebra deformations.

\bibitem{Wei2026}
S. Wei, Y. Zhai, Q. Lu, et al.,
The Riemann Hypothesis manifested in dynamical quantum phase transitions,
\emph{Nature Communications}, accepted 18 June 2026,
DOI 10.1038/s41467-026-74935-8.

\end{thebibliography}

\chapter{Reproduction}

From the source directory, run
\begin{lstlisting}[language=bash,numbers=none]
python projector_generated_special_functions_vii.py
PYTHONPATH=. pytest -q
pdflatex -interaction=nonstopmode Projector_Generated_Special_Functions_VII.tex
pdflatex -interaction=nonstopmode Projector_Generated_Special_Functions_VII.tex
\end{lstlisting}

\end{document}
